\section*{Capítulo \ref{ch:g11:vision-and-brain} --- La visión y el encéfalo}

\begin{solution}{exo:g11:vision-and-brain:1}
La luz del campo izquierdo cae sobre la mitad derecha de cada retina;
las fibras de la mitad derecha del \omterm{def:g11:the-eye:parts}{ojo} izquierdo cruzan en el quiasma y
las del \omterm{def:g11:the-eye:parts}{ojo} derecho se quedan a la derecha; juntas forman el tracto
óptico derecho, hacen relevo en el tálamo y llegan a la \omterm{prop:g11:vision-and-brain:pathway}{corteza visual
primaria} derecha.
\end{solution}

\begin{solution}{exo:g11:vision-and-brain:2}
Los dos nervios ópticos se encuentran; las fibras de la mitad nasal de
cada retina cruzan al otro lado, y las de la mitad temporal no. Más
allá, cada tracto lleva una mitad del campo visual.
\end{solution}

\begin{solution}{exo:g11:vision-and-brain:3}
El área del color (mundo visto en gris), el área del movimiento (sin
percepción del movimiento) y el área de las caras (caras no
reconocidas).
\end{solution}

\begin{solution}{exo:g11:vision-and-brain:4}
Plasticidad: la capacidad de las conexiones neuronales de reforzarse,
debilitarse, formarse o podarse según el uso. Periodo crítico: la etapa
temprana de la vida durante la cual la experiencia modela la corteza de
forma decisiva e irreversible.
\end{solution}

\begin{solution}{exo:g11:vision-and-brain:5}
Se une a los receptores de serotonina de las neuronas de las áreas
visuales, perturba el paso de las señales y produce alucinaciones sin
ninguna luz que les corresponda.
\end{solution}

\begin{solution}{exo:g11:vision-and-brain:6}
Después del quiasma, a la izquierda: el tracto óptico izquierdo, el
relevo talámico izquierdo o la \omterm{prop:g11:vision-and-brain:pathway}{corteza visual primaria} izquierda.
\end{solution}

\begin{solution}{exo:g11:vision-and-brain:7}
Las fibras que cruzan vienen de las mitades nasales de las dos retinas,
que ven las mitades externas (temporales) del campo: el paciente pierde
la mitad izquierda del campo del \omterm{def:g11:the-eye:parts}{ojo} izquierdo y la mitad derecha del
campo del \omterm{def:g11:the-eye:parts}{ojo} derecho --- una visión en túnel por los dos lados.
\end{solution}

\begin{solution}{exo:g11:vision-and-brain:8}
Normal: $12 + 20 + 36 + 20 = 88\%$. Privado: $2 + 3 + 5 + 30 = 40\%$,
de las que solo un 10\% de forma principal.
\end{solution}

\begin{solution}{exo:g11:vision-and-brain:9}
Tapar obliga a la corteza a usar la entrada del \omterm{def:g11:the-eye:parts}{ojo} débil, de modo que
sus conexiones se conservan y se refuerzan durante el periodo crítico
en lugar de podarse. En un adulto el periodo ha terminado: las
conexiones ya no se reorganizan a esa escala, así que el \omterm{def:g11:the-eye:parts}{ojo} débil no
puede recuperar su territorio.
\end{solution}

\begin{solution}{exo:g11:vision-and-brain:10}
Sus \omterm{def:g11:the-eye:photoreceptors}{conos} y sus \omterm{def:g11:the-eye:photoreceptors}{opsinas} son normales; lo que se ha perdido es la etapa
cortical que interpreta sus proporciones. Un registro eléctrico de la
respuesta de la retina a luces de colores sería normal en ella y
anormal en una persona daltónica; y su déficit es total y adquirido,
mientras que el de ellas es parcial e innato.
\end{solution}

\begin{solution}{exo:g11:vision-and-brain:11}
Años de práctica han reforzado y multiplicado las conexiones que
atienden a sus dedos, y han agrandado su territorio cortical. Sin
práctica, el área volvería a encogerse a lo largo de los años, como
muestran los estudios de los malabares.
\end{solution}

\begin{solution}{exo:g11:vision-and-brain:12}
Al principio ve luz, colores y manchas en movimiento, pero no reconoce
formas, ni caras, ni profundidad: la retina funciona, pero la corteza
no construyó, durante el periodo crítico, las conexiones que
interpretan sus señales. Aprenderá a usar la visión en parte y
despacio, pero nunca con la soltura de quien vio desde que nació.
\end{solution}

\begin{solution}{exo:g11:vision-and-brain:13}
En la fóvea, cada cono tiene su propia \omterm{def:g10:cells-common-unit:cell}{célula} ganglionar, así que el
centro del campo envía muchas más fibras por grado que la periferia,
donde cien \omterm{def:g11:the-eye:photoreceptors}{bastones} comparten una; la corteza reparte superficie en
proporción a las fibras y, por tanto, al detalle, y no a los grados de
campo.
\end{solution}

\begin{solution}{exo:g11:vision-and-brain:14}
La actividad de esa área es la percepción del color, venga la señal del
\omterm{def:g11:the-eye:parts}{ojo} o de la memoria; la percepción es lo que hacen las áreas, no lo que
entrega el \omterm{def:g11:the-eye:parts}{ojo}. Una alucinación es esa actividad puesta en marcha por
una droga en lugar de por la intención o por la luz.
\end{solution}

\begin{solution}{exo:g11:vision-and-brain:15}
El \omterm{def:g11:the-eye:parts}{ojo} convierte la luz en señales; ver lo hace el encéfalo. Un
paciente con los \omterm{def:g11:the-eye:parts}{ojos} intactos puede perder el color, el movimiento o
las caras por una lesión cortical; un gatito con los \omterm{def:g11:the-eye:parts}{ojos} intactos se
queda funcionalmente ciego de uno de ellos si no se deja que la corteza
lo conecte; y una droga que no toca ni el \omterm{def:g11:the-eye:parts}{ojo} ni la estructura de la
corteza hace que una persona vea lo que no está ahí. El \omterm{def:g11:the-eye:parts}{ojo} suministra;
el encéfalo ve.
\end{solution}

\begin{solution}{pb:g11:vision-and-brain:1}
\textbf{1.} El \omterm{def:g11:the-eye:parts}{ojo} izquierdo o su nervio óptico, antes del quiasma:
solo están afectadas las fibras de un \omterm{def:g11:the-eye:parts}{ojo}.

\textbf{2.} El tracto óptico, el tálamo o la \omterm{prop:g11:vision-and-brain:pathway}{corteza visual} derechos:
un déficit compartido por los dos \omterm{def:g11:the-eye:parts}{ojos} para la misma mitad del campo
solo puede estar donde las fibras de los dos \omterm{def:g11:the-eye:parts}{ojos} para esa mitad se han
reunido, es decir después del quiasma.

\textbf{3.} El propio quiasma: quedan cortadas las fibras que cruzan,
las de las mitades nasales de las dos retinas, que ven las mitades
temporales del campo.

\textbf{4.} El área del movimiento, más allá de la corteza primaria.
Que los campos y la agudeza estén intactos muestra que la corteza
primaria y todo lo anterior funcionan.

\textbf{5.} B, C y D tienen las retinas intactas; la de A puede estarlo
o no. Se comprueba examinando la retina con un oftalmoscopio y
registrando su respuesta eléctrica a los destellos.

\textbf{6.} La corteza derecha recibe la mitad izquierda del campo.

\textbf{7.} Los $2^\circ$ centrales: \qty{1}{cm}; una banda de
$2^\circ$ a $40^\circ$: \qty{0.1}{cm}. Cociente 10.

\textbf{8.} La fóvea apiña sus \omterm{def:g11:the-eye:photoreceptors}{conos}, cada uno con su propia \omterm{def:g10:cells-common-unit:cell}{célula}
ganglionar, así que los grados centrales envían muchas más fibras que
la periferia; la corteza les da proporcionalmente más superficie.

\textbf{9.} Un punto ciego pequeño en el centro del campo izquierdo,
que es incapacitante para leer; más hacia delante, una zona ciega mayor
en la periferia del campo izquierdo, que se nota menos.

\textbf{10.} Moviendo los \omterm{def:g11:the-eye:parts}{ojos}, el paciente lleva cada palabra a una
parte de la retina cuya corteza está intacta, y usa el borde de la
lesión; despacio, pero es posible.

\textbf{11.} $20 + 36 + 20 = 76\%$.

\textbf{12.} $2 + 3 + 5 + 30 = 40\%$, y solo un 10\% con el \omterm{def:g11:the-eye:parts}{ojo}
izquierdo como entrada principal o equivalente.

\textbf{13.} Su entrada llega a la corteza, pero sus conexiones se
debilitaron y se podaron por falta de uso; las conexiones del \omterm{def:g11:the-eye:parts}{ojo}
derecho se expandieron por ese espacio.

\textbf{14.} Que la reorganización se limita a un periodo crítico de la
primera etapa de la vida; su equivalente humano es la ambliopía, que
solo se desarrolla en la infancia y no tiene tratamiento después.

\textbf{15.} Sin ningún \omterm{def:g11:the-eye:parts}{ojo} favorecido no hay ocupación: el histograma
se mantiene más o menos simétrico pero con menos neuronas que
respondan en total, y la visión del gatito es mala con los dos \omterm{def:g11:the-eye:parts}{ojos} ---
la privación sin competencia debilita a los dos.

\textbf{16.} El encéfalo suprime la imagen del \omterm{def:g11:the-eye:parts}{ojo} desviado para evitar
la visión doble; sin uso, sus conexiones se podan durante el periodo
crítico y el \omterm{def:g11:the-eye:parts}{ojo} queda débil para siempre. Las gafas corrigen la
óptica, pero el \omterm{def:g11:the-eye:parts}{ojo} no se usa, así que la poda sigue adelante.

\textbf{17.} Tapar el \omterm{def:g11:the-eye:parts}{ojo} bueno obliga a la corteza a usar las señales
del \omterm{def:g11:the-eye:parts}{ojo} débil, y así conserva y refuerza sus conexiones.

\textbf{18.} El periodo crítico termina hacia los siete años; después,
la corteza ya no reasigna territorio a esa escala, así que las
conexiones perdidas del \omterm{def:g11:the-eye:parts}{ojo} débil no pueden reconstruirse.

\textbf{19.} En el adulto las conexiones ya están establecidas y son
estables; la ocupación por el otro \omterm{def:g11:the-eye:parts}{ojo} es un fenómeno de la corteza en
desarrollo, durante el periodo crítico, y no de la corteza madura.

\textbf{20.} A: \omterm{def:g11:the-eye:parts}{ojo} o nervio óptico; C: quiasma; B: tracto derecho,
tálamo o corteza primaria; D: el área del movimiento, más allá de ella.
La plasticidad de la corteza, limitada a un periodo crítico, hace que
tapar el \omterm{def:g11:the-eye:parts}{ojo} funcione a los cuatro años y fracase a los quince.
\end{solution}
